\section{Wstęp}
\label{sec:wstep}

Kierowca w transporcie drogowym przez cały czas pracy korzysta z dokumentów
przewozowych. Najważniejszym z nich jest międzynarodowy list przewozowy CMR,
którego rolę i zawartość określa Konwencja o umowie międzynarodowego przewozu
drogowego towarów \cite{konwencja_cmr}. W przewozach krajowych i przy odbiorze
towaru z magazynu dochodzą do tego dokumenty wydania zewnętrznego (WZ),
faktury i inne potwierdzenia. Większość z nich nadal ma formę papierową:
kierowca otrzymuje je przy załadunku, zbiera pieczątki i podpisy przy rozładunku,
a potem dostarcza oryginały do biura firmy przewozowej.

Papierowy obieg dokumentów powoduje typowe problemy. Dokumenty giną lub niszczą
się w kabinie pojazdu, biuro poznaje ich treść dopiero po powrocie kierowcy,
a rozliczenie zlecenia i wystawienie faktury się opóźnia. Kierowca ma jednak
przy sobie telefon z aparatem, więc może dokument sfotografować i przesłać go
zaraz po otrzymaniu.

Niniejsza praca opisuje projekt i implementację \emph{Aplikacji kierowcy}:
responsywnej aplikacji webowej (PWA), w której kierowca fotografuje dokument
telefonem, kadruje i kompresuje zdjęcie, a następnie przechowuje je
w zaszyfrowanej postaci, z podziałem na typy, statusy i wersje.
Rozdział~\ref{sec:cel} określa cel projektu i wymagania, rozdział~\ref{sec:przeglad}
przedstawia istniejące rozwiązania, a rozdziały~\ref{sec:architektura}
i~\ref{sec:implementacja} opisują architekturę i implementację systemu.
Pracę zamyka podsumowanie z wnioskami (rozdział~\ref{sec:podsumowanie}).
